% Generated draft: requires mathematical and source review.
\documentclass[11pt]{article}
\usepackage[margin=25mm]{geometry}
\usepackage{fontspec}
\setmainfont{Latin Modern Roman}
\usepackage{amsmath,amssymb,mathtools,unicode-math}
\setmathfont{Latin Modern Math}
\usepackage{enumitem,needspace,xurl,hyperref}
\setlength{\parindent}{0pt}
\setlength{\parskip}{5pt}
\setlength{\emergencystretch}{3em}
\newcommand{\N}{\mathbb{N}}
\newcommand{\Z}{\mathbb{Z}}
\newcommand{\Q}{\mathbb{Q}}
\newcommand{\R}{\mathbb{R}}
\newcommand{\C}{\mathbb{C}}
\newcommand{\F}{\mathbb{F}}
\newcommand{\A}{\mathbb{A}}
\newcommand{\PP}{\mathbb{P}}
\newcommand{\E}{\mathbb{E}}
\newcommand{\eps}{\varepsilon}
\newcommand{\dd}{\,\mathrm{d}}
\begin{document}
\section*{UnsolvedMath ID: 30005583}
\textbf{Source problem code:} catalogue entry
\subsection{Definitions and mathematical statement}
On a smooth complex surface with complex structure \(J\), a Kähler metric \(g\) is a Riemannian metric compatible with \(J\) whose associated two-form \(\omega(u,v)=g(Ju,v)\) is closed. It is Kähler--Einstein if \(\operatorname{Ric}(g)=\lambda g\) for a real constant \(\lambda\). For a singular complex space \(X\), this equation can be imposed on its regular locus \(X_{\mathrm{reg}}\); the meaning of extension across the singular set requires additional hypotheses. A smooth Einstein four-manifold is a smooth real four-dimensional manifold \(M\) equipped with a Riemannian metric \(h\) satisfying \(\operatorname{Ric}(h)=\mu h\) for some real constant \(\mu\).

\textbf{Source question, with unresolved scope.} Can every singular Kähler--Einstein metric be desingularized by smooth Einstein four-manifolds, allowing the desingularizing Einstein metrics to be non-Kähler? Here ``not necessarily Kähler'' permits Kähler metrics as well; it does not require the approximating metrics to be non-Kähler.

The supplied literature assessment describes the target as Kähler--Einstein orbifolds. Under a real four-dimensional orbifold interpretation, the intended question is whether every such \((X,g,J)\) admits smooth Einstein four-manifolds \((M_i,h_i)\) that desingularize \((X,g)\). This is a conditional interpretation, not a fully specified conjecture: the source does not define the admissible singular spaces or the desingularization relation. In particular, it does not specify a convergence topology, identifications away from the singularities, completeness or compactness requirements, boundary conditions, permitted changes of topology, or the relation between the Einstein constants of \(g\) and \(h_i\). The original question does not explicitly restrict the singular metric to real dimension four, although the proposed desingularizing manifolds have that dimension. These choices must be supplied before the universal assertion has a precise mathematical meaning.
\subsection{Short English statement}
The question asks whether every singular Kähler--Einstein metric can arise through desingularization by smooth Einstein four-manifolds, with no requirement that the smooth metrics be Kähler. The supplied assessment discusses orbifolds, but the admissible singularities and the meaning of desingularization remain unspecified. This is a statement-only draft requiring human mathematical review.
\subsection{Sources}
\begin{itemize}
\item Local UnsolvedMath problem registry, permanent ID \(30005583\), source problem code \texttt{catalogue entry}, title \emph{Desingularization of Kähler--Einstein Metrics}, in the Oberwolfach Reports Open Problems collection. The supplied record labels the problem \texttt{open} and the statement \texttt{corrected\_verified}; these are supplied registry labels, not independent verification here.
\item Supplied report citation: \emph{Differentialgeometrie im Grossen (2024)}, Oberwolfach Reports, DOI \texttt{10.4171/owr/2023/29}, pp.\ 1617--1670. \url{https://doi.org/10.4171/owr/2023/29}. The record gives proposed year 2023. The relationship between that year, the DOI's 2023 component, and the citation's 2024 date is not clarified in the supplied material.
\item The original extract associates the question with \emph{Holonomy of limits of Einstein 4-manifolds}, Tristan Ozuch, joint work with Claude LeBrun. No separate proposer is identified. The extract includes the reference label \texttt{[And10]}, but supplies no bibliographic entry for it.
\item Supplied literature source labeled ``Ozuch'': \url{https://doi.org/10.1002/cpa.22129}.
\item Supplied literature source labeled ``LeBrun--Ozuch'': \url{https://arxiv.org/abs/2601.19215}.
\item Literature assessment dated exactly \texttt{2026-08-22}: the supplied curation labels the question open and states that general desingularization of Kähler--Einstein orbifolds by Einstein four-manifolds is not settled; known obstruction results and conditional Hermitian or anti-self-dual gluing theorems do not establish the universal assertion. This is supplied dated curation, not an independently verified literature assessment. The supplied background explicitly requests independent verification of time-sensitive status and cited claims.
\item The registry reports that the official Oberwolfach report was checked directly and that the cleaned statement preserves the mathematical target while removing extraction artifacts or source-dependent wording. That verification claim is part of the supplied provenance; no source has been independently visited or checked for this draft.
\end{itemize}
\end{document}
